\chapter*{Résumé}
\addcontentsline{toc}{chapter}{Résumé}

Dans un contexte de transformation digitale du secteur bancaire, la donnée constitue un levier stratégique essentiel pour la prise de décision, la gestion des risques et l’amélioration de la performance globale. Cependant, la qualité des données représente un enjeu majeur, car des données incomplètes, incohérentes ou erronées peuvent impacter significativement les processus métiers et les analyses.

Ce projet s’inscrit dans le cadre des initiatives de Data Governance au sein de CIH Bank. Il vise à concevoir et mettre en œuvre un pipeline automatisé de contrôle de la qualité des données, permettant d’assurer la fiabilité, la traçabilité et la cohérence des informations exploitées.

La solution proposée repose sur le traitement et le contrôle des données du domaine des Tiers (données KYC et CSP) ingérées dans le Data Lake (sous forme de tables Hive/ORC), en intégrant des règles de contrôle couvrant plusieurs dimensions de la qualité telles que la complétude, la validité, la cohérence et l’unicité.

L’automatisation des contrôles permet d’instaurer une surveillance continue de la qualité des données, avec la production d’indicateurs (KPI) facilitant le pilotage et la prise de décision. Ce projet contribue ainsi à renforcer la gouvernance des données et à soutenir les initiatives analytiques et décisionnelles de la banque.
